% GENERATED FILE. Edit canonical lessons, then run npm run generate:books.
% canonical-source-hash: fnv1a64:1975841a28faad0d
% canonical-lessons: KA-C378-drstikona, KA-C378-niluvu, KA-C378-tatastha, KA-C378-paksapata, KA-C378-pratipadisu

\chapter{Taking a Stand}
\label{ch:ka-taking-a-stand}
\hlchaptermodality{378}

\begin{chapteropening}
\textbf{By the end of this chapter:} \emph{I can say a point of view, a stance, neutral, bias and to advocate.}

Complete the last lesson of chapter 378: I can say a point of view, a stance, neutral, bias and to advocate.
\end{chapteropening}

\section[\texorpdfstring{\kn{ದೃಷ್ಟಿಕೋನ} (dṛṣṭikōna)}{ದೃಷ್ಟಿಕೋನ (dṛṣṭikōna)}]{\kn{ದೃಷ್ಟಿಕೋನ} (dṛṣṭikōna) — a point of view}

\label{lesson:KA-C378-drstikona}



\begin{quote}
Before the new one: say the Kannada for to object, then the Kannada for reasonable.
\end{quote}

\subsection*{You'll want to know: \kn{ದೃಷ್ಟಿಕೋನ}}
\textbf{\kn{ದೃಷ್ಟಿಕೋನ}} — \emph{dṛṣṭikōna} — \textquotedblleft{}a point of view\textquotedblright{}.

\begin{quote}
\kn{ಪ್ರತಿಯೊಬ್ಬರಿಗೂ} \kn{ಬೇರೆ} \kn{ಬೇರೆ} \kn{ದೃಷ್ಟಿಕೋನ} \kn{ಇರುತ್ತದೆ}.
\end{quote}

\begin{quote}
\emph{pratiyobbarigū bēre bēre dṛṣṭikōna iruttade.}
\end{quote}

\begin{quote}
Everyone has a different point of view.
\end{quote}

Word by word it is sight + angle.

\begin{cousinweb}
\textbf{\kn{ದೃಷ್ಟಿಕೋನ}} (\emph{dṛṣṭikōna}) is built from Sanskrit \emph{dṛṣṭi}, \textquotedblleft{}sight, view\textquotedblright{}, and \emph{koṇa}, \textquotedblleft{}a corner, an angle\textquotedblright{} (Monier-Williams): the angle you look from.
\end{cousinweb}

\begin{grammarlens}[title={What you've built}]
One more word for giving reasons and opinions.
\end{grammarlens}

\subsection*{Your turn}
\begin{itemize}
\raggedright
  \item \emph{Say it:} \emph{dṛṣṭikōna}
  \item \emph{Say it:} \emph{dṛṣṭikōna}, once more
  \item \emph{Say it:} \emph{pratiyobbarigū bēre bēre dṛṣṭikōna iruttade.} — Everyone has a different point of view.
  \item \emph{Recall:} say the Kannada for to object, then the Kannada for reasonable, then say \emph{dṛṣṭikōna} again
\end{itemize}

\begin{tcolorbox}[breakable,colback=teal!4,colframe=teal!35!black,title={Before you move on}]
What does \kn{ದೃಷ್ಟಿಕೋನ} mean? (\textquotedblleft{}A point of view\textquotedblright{}.) Say it once more.
\end{tcolorbox}
\section[\texorpdfstring{\kn{ನಿಲುವು} (niluvu)}{ನಿಲುವು (niluvu)}]{\kn{ನಿಲುವು} (niluvu) — a stance; a position}

\label{lesson:KA-C378-niluvu}



\begin{quote}
Before the new one: say the Kannada for reasonable, then the Kannada for a point of view.
\end{quote}

\subsection*{You'll want to know: \kn{ನಿಲುವು}}
\textbf{\kn{ನಿಲುವು}} — \emph{niluvu} — \textquotedblleft{}a stance; a position\textquotedblright{}.

\begin{quote}
\kn{ಈ} \kn{ವಿಷಯದಲ್ಲಿ} \kn{ನಿಮ್ಮ} \kn{ನಿಲುವು} \kn{ಏನು}?
\end{quote}

\begin{quote}
\emph{ī viṣayadalli nimma niluvu ēnu?}
\end{quote}

\begin{quote}
What is your position on this matter?
\end{quote}

From \kn{ನಿಲ್ಲು}, to stand: where one stands.

\begin{cousinweb}
\textbf{\kn{ನಿಲುವು}} (\emph{niluvu}) is native Dravidian, a noun from \emph{nil}, \textquotedblleft{}to stand\textquotedblright{} (DEDR 3675), the root of \kn{ನಿಲ್ಲು}.
\end{cousinweb}

\begin{grammarlens}[title={What you've built}]
One more word for giving reasons and opinions.
\end{grammarlens}

\subsection*{Your turn}
\begin{itemize}
\raggedright
  \item \emph{Say it:} \emph{niluvu}
  \item \emph{Say it:} \emph{niluvu}, once more
  \item \emph{Say it:} \emph{ī viṣayadalli nimma niluvu ēnu?} — What is your position on this matter?
  \item \emph{Recall:} say the Kannada for reasonable, then the Kannada for a point of view, then say \emph{niluvu} again
\end{itemize}

\begin{tcolorbox}[breakable,colback=teal!4,colframe=teal!35!black,title={Before you move on}]
What does \kn{ನಿಲುವು} mean? (\textquotedblleft{}A stance; a position\textquotedblright{}.) Say it once more.
\end{tcolorbox}
\section[\texorpdfstring{\kn{ತಟಸ್ಥ} (taṭastha)}{ತಟಸ್ಥ (taṭastha)}]{\kn{ತಟಸ್ಥ} (taṭastha) — neutral}

\label{lesson:KA-C378-tatastha}



\begin{quote}
Before the new one: say the Kannada for a point of view, then the Kannada for a stance.
\end{quote}

\subsection*{You'll want to know: \kn{ತಟಸ್ಥ}}
\textbf{\kn{ತಟಸ್ಥ}} — \emph{taṭastha} — \textquotedblleft{}neutral\textquotedblright{}.

\begin{quote}
\kn{ಜಗಳದಲ್ಲಿ} \kn{ನಾನು} \kn{ತಟಸ್ಥವಾಗಿ} \kn{ಇದ್ದೆ}.
\end{quote}

\begin{quote}
\emph{jagaḷadalli nānu taṭasthavāgi idde.}
\end{quote}

\begin{quote}
In the quarrel I stayed neutral.
\end{quote}

Word by word it is standing on the bank, out of the stream.

\begin{cousinweb}
\textbf{\kn{ತಟಸ್ಥ}} (\emph{taṭastha}) is borrowed from Sanskrit \emph{taṭastha}, \textquotedblleft{}standing on a bank\textquotedblright{} (Monier-Williams), from \emph{taṭa}, \textquotedblleft{}a bank\textquotedblright{}, and \emph{stha}, \textquotedblleft{}standing\textquotedblright{}.
\end{cousinweb}

\begin{grammarlens}[title={What you've built}]
One more word for giving reasons and opinions.
\end{grammarlens}

\subsection*{Your turn}
\begin{itemize}
\raggedright
  \item \emph{Say it:} \emph{taṭastha}
  \item \emph{Say it:} \emph{taṭastha}, once more
  \item \emph{Say it:} \emph{jagaḷadalli nānu taṭasthavāgi idde.} — In the quarrel I stayed neutral.
  \item \emph{Recall:} say the Kannada for a point of view, then the Kannada for a stance, then say \emph{taṭastha} again
\end{itemize}

\begin{tcolorbox}[breakable,colback=teal!4,colframe=teal!35!black,title={Before you move on}]
What does \kn{ತಟಸ್ಥ} mean? (\textquotedblleft{}Neutral\textquotedblright{}.) Say it once more.
\end{tcolorbox}
\section[\texorpdfstring{\kn{ಪಕ್ಷಪಾತ} (pakṣapāta)}{ಪಕ್ಷಪಾತ (pakṣapāta)}]{\kn{ಪಕ್ಷಪಾತ} (pakṣapāta) — bias; partiality}

\label{lesson:KA-C378-paksapata}



\begin{quote}
Before the new one: say the Kannada for a stance, then the Kannada for neutral.
\end{quote}

\subsection*{You'll want to know: \kn{ಪಕ್ಷಪಾತ}}
\textbf{\kn{ಪಕ್ಷಪಾತ}} — \emph{pakṣapāta} — \textquotedblleft{}bias; partiality\textquotedblright{}.

\begin{quote}
\kn{ಶಿಕ್ಷಕರು} \kn{ಯಾರಿಗೂ} \kn{ಪಕ್ಷಪಾತ} \kn{ಮಾಡುವುದಿಲ್ಲ}.
\end{quote}

\begin{quote}
\emph{śikṣakaru yārigū pakṣapāta māḍuvudilla.}
\end{quote}

\begin{quote}
The teacher shows no bias to anyone.
\end{quote}

\kn{ಪಕ್ಷಪಾತ} \kn{ಮಾಡು}, to be partial.

\begin{cousinweb}
\textbf{\kn{ಪಕ್ಷಪಾತ}} (\emph{pakṣapāta}) is borrowed from Sanskrit \emph{pakṣapāta}, \textquotedblleft{}siding with, partiality\textquotedblright{} (Monier-Williams), from \emph{pakṣa}, \textquotedblleft{}a side, a wing\textquotedblright{}, and \emph{pāta}, \textquotedblleft{}falling\textquotedblright{}.
\end{cousinweb}

\begin{grammarlens}[title={What you've built}]
One more word for giving reasons and opinions.
\end{grammarlens}

\subsection*{Your turn}
\begin{itemize}
\raggedright
  \item \emph{Say it:} \emph{pakṣapāta}
  \item \emph{Say it:} \emph{pakṣapāta}, once more
  \item \emph{Say it:} \emph{śikṣakaru yārigū pakṣapāta māḍuvudilla.} — The teacher shows no bias to anyone.
  \item \emph{Recall:} say the Kannada for a stance, then the Kannada for neutral, then say \emph{pakṣapāta} again
\end{itemize}

\begin{tcolorbox}[breakable,colback=teal!4,colframe=teal!35!black,title={Before you move on}]
What does \kn{ಪಕ್ಷಪಾತ} mean? (\textquotedblleft{}Bias; partiality\textquotedblright{}.) Say it once more.
\end{tcolorbox}
\section[\texorpdfstring{\kn{ಪ್ರತಿಪಾದಿಸು} (pratipādisu)}{ಪ್ರತಿಪಾದಿಸು (pratipādisu)}]{\kn{ಪ್ರತಿಪಾದಿಸು} (pratipādisu) — to advocate; to argue for}

\label{lesson:KA-C378-pratipadisu}



\begin{quote}
Before the new one: say the Kannada for neutral, then the Kannada for bias.
\end{quote}

\subsection*{You'll want to know: \kn{ಪ್ರತಿಪಾದಿಸು}}
\textbf{\kn{ಪ್ರತಿಪಾದಿಸು}} — \emph{pratipādisu} — \textquotedblleft{}to advocate; to argue for\textquotedblright{}.

\begin{quote}
\kn{ಅವರು} \kn{ಎಲ್ಲರಿಗೂ} \kn{ಉಚಿತ} \kn{ಶಿಕ್ಷಣವನ್ನು} \kn{ಪ್ರತಿಪಾದಿಸಿದರು}.
\end{quote}

\begin{quote}
\emph{avaru ellarigū ucita śikṣaṇavannu pratipādisidaru.}
\end{quote}

\begin{quote}
They argued for free education for all.
\end{quote}

Past (I …): \kn{ಪ್ರತಿಪಾದಿಸಿದೆ}, \emph{pratipādiside}.

What you argue for takes -\kn{ಅನ್ನು}.

\begin{cousinweb}
\textbf{\kn{ಪ್ರತಿಪಾದಿಸು}} (\emph{pratipādisu}) is built on Sanskrit \emph{pratipādana}, \textquotedblleft{}presenting, setting forth\textquotedblright{} (Monier-Williams), with the Kannada verb ending -\emph{isu}.
\end{cousinweb}

\begin{grammarlens}[title={What you've built}]
One more word for giving reasons and opinions.
\end{grammarlens}

\subsection*{Your turn}
\begin{itemize}
\raggedright
  \item \emph{Say it:} \emph{pratipādisu}
  \item \emph{Say it:} \emph{pratipādisu}, once more
  \item \emph{Say it:} \emph{avaru ellarigū ucita śikṣaṇavannu pratipādisidaru.} — They argued for free education for all.
  \item \emph{Recall:} say the Kannada for neutral, then the Kannada for bias, then say \emph{pratipādisu} again
\end{itemize}

\begin{tcolorbox}[breakable,colback=teal!4,colframe=teal!35!black,title={Before you move on}]
What does \kn{ಪ್ರತಿಪಾದಿಸು} mean? (\textquotedblleft{}To advocate; to argue for\textquotedblright{}.) Say it once more.
\end{tcolorbox}
